\section{Gaussian binary least squares}
\label{sec:binary-least-squares}

Binary least squares separates finding an optimal vertex from certifying its
value with a particular relaxation.  In the Gaussian model considered here,
the planted vertex is the unique integer optimum above a logarithmic signal
level.  The box relaxation can nevertheless require a superpolynomial
certificate.  At a linear signal level, a sufficient condition for a linear
variable-branching tree has a sharp threshold.  A different, logarithmic
threshold determines when the root box relaxation agrees with its
nonnegative least-squares relaxation.  These statements concern different
events and different certificate classes.

\subsection{Model and certificate conventions}
\label{bls:model}

Let $M=M_N\ge N$, let $H\in\R^{M\times N}$ have independent standard
Gaussian entries, and let $w\sim N(0,I_M)$ be independent of $H$.  For a
deterministic $x^*\in\{-1,1\}^N$ and $\rho>0$, set
\begin{equation}
 A=\sqrt{\rho/N}\,H,\qquad y=Ax^*+w,\qquad
 f(x)=\|y-Ax\|^2,\qquad
 \OPT=\min_{x\in\{-1,1\}^N}f(x),\qquad W=\|w\|^2.
 \label{bls:model-equation}
\end{equation}
Thus $W=f(x^*)$ and $\rho$ is normalized per observation.  Write
$\beta_N=M/N$; asymptotic statements below assume
$\beta_N\to\beta\in[1,\infty)$.  All probabilities are uniform in the
deterministic planted signs.  Dividing the residual by a positive noise
standard deviation $\sigma$ replaces $\rho$ by $\rho/\sigma^2$.

The node oracle minimizes $f$ over the node's convex set inside
$[-1,1]^N$.  In particular, a variable-branching node minimizes $f$ over
the box obtained by fixing some coordinates to binary values.  No identity
$x_i^2=1$ is imposed on free coordinates.  Diagonal quadratic
reformulations, product lifts, and semidefinite relaxations change this
oracle and are outside the certificate lower bound proved here.

For a target $\tau$, define the class number
\begin{equation}
 \kappa_\tau=\min\left\{q:
   \{-1,1\}^N=\bigcup_{j=1}^q I_j,\quad
   \min_{x\in\operatorname{conv} I_j}f(x)\ge\tau
   \text{ for every }j\right\}.
 \label{bls:class-number}
\end{equation}
Empty classes can be discarded.  This is the minimum number of leaves of
a certificate allowing arbitrary convex pieces that cover the binary
vertices: a leaf containing $I_j$ also contains its convex hull, and the
convex hulls themselves give a certificate.  Every variable or general
split tree using this oracle has at least $\kappa_\tau$ leaves.  Adding
convex cuts valid for all the binary vertices assigned to a leaf preserves this
argument, because those cuts contain their convex hull.  Incumbent-based
reductions also need certificates for the removed vertices; their cost is
handled by the augmented accounting of the integer framework, rather than
by treating the retained leaves alone as a cover.

Absorb the planted signs into the columns by writing
$b_i=x_i^*a_i$, $B=A\diag(x^*)$, and
$u_i=1-x_i^*x_i$.  Then
\begin{equation}
 f(x)=F(u):=\|w+Bu\|^2,\qquad u\in[0,2]^N.
 \label{bls:error-coordinates}
\end{equation}
The planted vertex is $u=0$, and flipping the coordinates in $S$ gives
$u=2\mathbf 1_S$.  If a node fixes the wrong values on $T$ and leaves the
coordinates $J$ free, its box value is
\begin{equation}
 r(J,T)=\min_{u\in[0,2]^J}\|v_T+B_Ju\|^2,
 \qquad v_T=w+2B_T\mathbf 1.
 \label{bls:node-equation}
\end{equation}
The residual $v_T$ is independent of the free matrix $B_J$.  Shared noise
creates dependence between different nodes, but does not change this
one-node independence.

\subsection{Recovery and root exactness are different events}
\label{bls:recovery-root}

We first record the recovery range used when interpreting a bound relative
to the planted value as a bound relative to $\OPT$.

\begin{proposition}[Planted recovery and a value comparison]
\label{bls:ml}
For every fixed $\delta>0$, if
$\beta_N\rho_N\ge(2+\delta)\log N$, then $x^*$ is the unique integer
minimizer with probability tending to one.  For every deterministic
positive sequence $\rho_N$, with probability tending to one,
\begin{equation}
 0\le W-\OPT\le
 4N\log\bigl(1+e^{-0.14\beta_N\rho_N}\bigr)+4\log N.
 \label{bls:ml-value-gap}
\end{equation}
\end{proposition}

The proof is a first-moment calculation over flipped sets, given in
\Cref{app:bls-ml-proof}.  At $\rho_N=\theta N$ with fixed $\theta>0$,
it in particular proves unique planted optimality without a box-decoder
theorem.  For square systems, the recovery comparison at a logarithmic
signal level is consistent with the sufficient condition in
\citet{hassibiHansenDimakisAlshamaryXu2014OptimizedMCMC}; we use the fixed
margin $(2+\delta)\log N$.  Recovery is not asserted for every sequence
of order $\log N$.

Even when $x^*$ is the integer optimum, it usually does not minimize the
root box relaxation.  Let $R=\min_{x\in[-1,1]^N}f(x)$.

\begin{proposition}[Two root exactness events]
\label{bls:root-exactness}
For every $M\ge N$ and $\rho>0$,
\begin{equation}
 \Pbb\{x^*\text{ minimizes }f\text{ on }[-1,1]^N\}=2^{-N}.
 \label{bls:planted-root-probability}
\end{equation}
The root minimizer is unique almost surely, and $R=\OPT$ if and only if
that minimizer is a vertex.  Moreover,
\begin{align}
 \Pbb\{R=\OPT\}
 &\le\sum_{k=0}^N\binom Nk2^{k-N}
             (1+4\rho k/N)^{-M/2}\notag\\
 &\le\left[\frac{1+2(1+4\rho)^{-\beta_N/2}}2\right]^N.
 \label{bls:global-root-probability}
\end{align}
\end{proposition}

The last bound decreases exponentially when
$(1+4\rho)^{\beta_N/2}>2$ with a fixed margin.  The exact planted
probability and the global vertex probability are distinct: at low signal
levels another vertex can be the integer optimum.  Rounding the box
minimizer is a further event.  In particular, a successful rounded decoder
does not imply $R=\OPT$.  The proof in
\Cref{app:bls-root-exactness-proof} uses only the box optimality conditions
and a Gaussian exponential moment.

\subsection{Finite Gaussian nonnegative least squares}
\label{bls:nnls}

The next results supply the one-node laws and the tail estimates needed
for simultaneous conclusions.  They also isolate the role of selecting a
support from the data.  The binomial face-dimension law and the chi-square
projection mixture are classical Gaussian cone facts; see
\citet{hugSchneider2016RandomConicalTessellations} and the master Steiner
formula of \citet{mccoyTropp2014SteinerFormulas}.  We give a direct proof
because its invariant-event identity also transfers coefficient laws.
The Gaussian regression Student density itself is classical.  The finite
coefficient bound below is the ingredient used for the branch-and-bound
applications.

\begin{lemma}[An invariant-event identity]
\label{bls:sign-orbit}
Let $C\in\R^{M\times n}$, $n\le M$, have full column rank almost surely,
and let $v$ be independent of $C$.  Assume the joint column law of $C$ is
invariant under every independent column sign change.  For each nonempty
$S\subseteq[n]$, put
\[
 d^S=-(C_S^TC_S)^{-1}C_S^Tv,\qquad
 L_S=\operatorname{span}(C_S),\qquad z^S=(I-P_{L_S})v.
\]
For the empty support set $L_\varnothing=\{0\}$ and $z^\varnothing=v$.
Assume that almost surely every coordinate of every $d^S$ is nonzero,
and $c_j^Tz^S\ne0$ for every $S\subseteq[n]$ and $j\notin S$.
Let $u^\circ$ minimize $\|v+Cu\|^2$ over $u\ge0$, and let
$\mathcal S=\{j:u_j^\circ>0\}$.  If $E_S$ is any measurable event
invariant under every column sign change, then
\begin{equation}
 \Pbb\{\mathcal S=S,\ E_S\}=2^{-n}\Pbb(E_S).
 \label{bls:invariant-identity}
\end{equation}
In particular, $\mathcal S$ is uniform over all subsets of $[n]$.
\end{lemma}

\begin{proof}
For fixed $C,v$ satisfying the hypotheses, the support is $S$ precisely
when $d_j^S>0$ for $j\in S$ and $c_j^Tz^S>0$ for $j\notin S$.
These are the necessary and sufficient first-order conditions for the
strictly convex problem.  A column sign change preserves $L_S,z^S$,
changes the sign of the corresponding coordinate of $d^S$ on $S$, and
changes the sign of the corresponding residual inner product off $S$.
Exactly one of the $2^n$ sign patterns satisfies all these conditions.
Multiplying this pointwise sign count by the invariant indicator of $E_S$
and taking expectations proves \eqref{bls:invariant-identity}.  Conditions
indexed by an empty set are vacuous, so the argument includes the empty
and full supports.
\end{proof}

The explicit nonzero conditions matter: requiring only that $v$ avoid
column spans would not suffice.  For independent Gaussian columns and
independent isotropic Gaussian $v$ of positive variance, the stated
conditions hold automatically.  Conditional on a full-rank matrix, each
coefficient and each off-support residual inner product is a nonzero
linear functional of $v$.  There are finitely many such functionals.

\begin{theorem}[Gaussian NNLS coefficient and projection laws]
\label{bls:nnls-law}
Let $C=aH$, where $a>0$ and $H\in\R^{M\times n}$ has independent
standard Gaussian entries, $n\le M$.  Let $v\sim N(0,c^2I_M)$ be
independent of $H$, with $c>0$, and define $u^\circ,\mathcal S$ as in
\Cref{bls:sign-orbit}.  Then $K=|\mathcal S|\sim\operatorname{Bin}(n,1/2)$.
For a specified support $S$ of size $s\ge1$, conditional on
$\mathcal S=S$, the positive coefficient vector has the law
\begin{equation}
 (u_j^\circ)_{j\in S}\ \stackrel d=
 \frac ca\frac{|Z_s|}{\sqrt Q},\qquad
 Z_s\sim N(0,I_s),\quad Q\sim\chi^2_{M-s+1},
 \label{bls:coefficient-mixture}
\end{equation}
where $Z_s,Q$ are independent and the absolute value is coordinatewise.
Equivalently, its unconditioned signed regression vector $d^S$ has density
\begin{equation}
 p_s(d)=\left(\frac ac\right)^s
 \frac{\Gamma((M+1)/2)}{\pi^{s/2}\Gamma((M-s+1)/2)}
 \left(1+\frac{a^2}{c^2}\|d\|^2\right)^{-(M+1)/2}.
 \label{bls:student-density}
\end{equation}
Writing $\overline\Phi$ for the standard Gaussian upper-tail function,
for $t\ge0$ the maximum has the exact distribution
\begin{equation}
 \Pbb\{\max_j u_j^\circ\le t\}
 =2^{-n}\sum_{s=0}^n\binom ns
  \E_{Q\sim\chi^2_{M-s+1}}
  \left[1-2\overline\Phi\!\left(\frac{at\sqrt Q}{c}\right)\right]^s,
 \label{bls:maximum-cdf}
\end{equation}
where the $s=0$ summand is one.  In particular, the atom at zero is $2^{-n}$.

If $X=\|Cu^\circ\|^2$ and $Z=\|v+Cu^\circ\|^2$, then conditional on
$K=s$, $X/c^2$ and $Z/c^2$ are independent chi-square variables with
$s$ and $M-s$ degrees of freedom.  Thus
$Y=X/(X+Z)$ has the corresponding beta mixture, with point masses at the
endpoints when a degree is zero.  For $0\le t\le1$,
\begin{equation}
 \Pbb\{|Y-n/(2M)|>2t\}\le6e^{-Mt^2/64}.
 \label{bls:projection-concentration}
\end{equation}
\end{theorem}

The auxiliary chi-square denominator in \eqref{bls:coefficient-mixture}
is common to all coordinates on that support.  It is not a set of
independent one-coordinate denominators, nor a selected singular value.
The representations concern one NNLS problem and assert no independence
between different branch-and-bound nodes.  The complete density and
projection proofs are in \Cref{app:bls-nnls-proofs}.

\begin{theorem}[A finite maximum-coefficient tail]
\label{bls:nnls-tail}
Under the hypotheses of \Cref{bls:nnls-law}, for every $t>0$ set
$q=(1+a^2t^2/c^2)^{-1/2}$.  Then
\begin{equation}
 \Pbb\{\max_j u_j^\circ>t\}
 \le\frac n2\,q^{M-n+1}\left(\frac{1+q}{2}\right)^{n-1}.
 \label{bls:maximum-tail}
\end{equation}
\end{theorem}

\begin{proof}
On support $S$, the positive coordinates are $d^S$.  The event
$\max_{j\in S}|d_j^S|>t$ is invariant under column sign changes, so
\Cref{bls:sign-orbit} and a coordinate union bound give
\[
 \Pbb\{\max_j u_j^\circ>t\}
 \le2^{-n}\sum_{S\ne\varnothing}\sum_{j\in S}
                   \Pbb\{|d_j^S|>t\}.
\]
For a fixed support of size $s$ and $j\in S$, let $P$ project orthogonally
to the other $s-1$ columns of $H_S$.  The regression identity gives
$d_j^S=-h_j^TPv/(aQ_j)$, where $Q_j=\|Ph_j\|^2\sim\chi^2_{M-s+1}$.
Conditional on $H_S$, this coefficient is
$N(0,c^2/(a^2Q_j))$.  The elementary bound
$\Pbb\{|G|>z\}\le e^{-z^2/2}$ for a standard normal $G$ and the
chi-square Laplace transform therefore give
\[
 \Pbb\{|d_j^S|>t\}
 \le\E e^{-a^2t^2Q_j/(2c^2)}=q^{M-s+1}.
\]
Consequently the required sum is at most
\[
 2^{-n}\sum_{s=1}^n\binom ns s q^{M-s+1}
 =\frac n2 q^{M-n+1}\left(\frac{1+q}{2}\right)^{n-1}.
\]
This argument includes $s=M$, where the regression denominator has one
degree of freedom.
\end{proof}

\subsection{A sharp sufficient criterion for a linear variable tree}
\label{bls:linear-trees}

Let $r_i$ be the box bound at the node fixing $x_i=-x_i^*$ and leaving
all other coordinates free.  Criterion C1 at $x^*$ requires all these
bounds to reach the target.  Its usefulness is monotonicity: every node
off the path to $x^*$ is contained in one such single-wrong-fixing node.
The path lemma, \Cref{lattice:path}, then gives at most $2N+1$ nodes.
For weak bounds $r_i\ge\OPT$, the optimal incumbent must be available
when off-path nodes are checked.  For strict bounds $r_i>\OPT$, best-bound
search gives the same conclusion without an initial optimal incumbent.
The remaining singleton has exact bound $f(x^*)$, as required by the
path lemma.

\begin{theorem}[Simultaneous C1 threshold for square and tall systems]
\label{bls:c1-threshold}
Assume $\rho_N=\theta N$ with fixed $\theta>0$ and
$\beta_N\to\beta\ge1$.  With probability tending to one, the orthant
minimizer at every single-wrong-fixing node belongs to $[0,2]^{N-1}$.
Moreover,
\begin{equation}
 \max_{1\le i\le N}\left|
 \frac{r_i-W}{N}-\left[2\theta(2\beta-1)-\frac12\right]
 \right|\xrightarrow{\Pbb}0.
 \label{bls:uniform-c1-gap}
\end{equation}
Put $\theta_c=[4(2\beta-1)]^{-1}$.  For each fixed
$\epsilon\in(0,1)$, the planted vertex is uniquely optimal with
probability tending to one, and the following conclusions hold:
\begin{enumerate}[label=(\alph*)]
\item If $\theta\ge(1+\epsilon)\theta_c$, then all
$r_i\ge\OPT+\epsilon N/4$ with probability tending to one.  Every
variable-branching search satisfying the incumbent or best-bound
qualification above has at most $2N+1$ nodes.
\item If $\theta\le(1-\epsilon)\theta_c$, then all
$r_i\le\OPT-\epsilon N/4$ with probability tending to one.  Thus C1
fails at every single-wrong-fixing node.
\end{enumerate}
\end{theorem}

\begin{proof}
Here the signed columns are independent $N(0,\theta I_M)$.  At node $i$,
$v_i=w+2b_i\sim N(0,(1+4\theta)I_M)$ is independent of $B_{-i}$.
Apply \Cref{bls:nnls-tail} with
$n=N-1$, $a^2=\theta$, $c^2=1+4\theta$, and $t=2$.  Writing
\[
 q_\theta=\sqrt{\frac{1+4\theta}{1+8\theta}},\qquad
 b_\theta=\frac{1+q_\theta}{2}<1,
\]
a union bound over the $N$ nodes bounds the probability of any upper-box
violation by
\begin{equation}
 \frac{N(N-1)}2 q_\theta^{M-N+2}b_\theta^{N-2}.
 \label{bls:all-node-inactivity}
\end{equation}
This tends to zero exponentially, including when $M=N$.  No independence
of node events is used.

Define the orthant value first:
$\widetilde r_i=\min_{u\ge0}\|v_i+B_{-i}u\|^2$.
By \Cref{bls:nnls-law}, it has the marginal law
\[
 \widetilde r_i/(1+4\theta)\mid K_i=s\sim\chi^2_{M-s},
 \qquad K_i\sim\operatorname{Bin}(N-1,1/2).
\]
Binomial and chi-square exponential bounds, followed by a union bound over
$i$, give the following estimate.  Specifically, take
$s=\sqrt{2N\log N}$ in the binomial bounds and $x=4\log N$ in the
conditional chi-square bounds; their total failure probability is
$O(N^{-3})$.
\[
 \max_i\left|\widetilde r_i-
 (1+4\theta)\left(M-\frac{N-1}2\right)\right|
 =O_{\Pbb}(\sqrt{N\log N}).
\]
Also $W=M+O_{\Pbb}(\sqrt{N\log N})$.  On the event complementary to
\eqref{bls:all-node-inactivity}, strict convexity gives
$r_i=\widetilde r_i$ for every $i$.  Hence, uniformly over $i$,
\begin{equation}
 r_i-W=N\left[2\theta(2\beta_N-1)-\frac12\right]
       +\frac{1+4\theta}2+O_{\Pbb}(\sqrt{N\log N}),
 \label{bls:finite-c1-expansion}
\end{equation}
which proves \eqref{bls:uniform-c1-gap}.  The elementary Gaussian
concentration estimates used here are proved in
\Cref{app:bls-probability-lemmas}.

\Cref{bls:ml} gives $\OPT=W$ and unique planted optimality with probability
tending to one.  The limiting gap equals
$\tfrac12(\theta/\theta_c-1)$, so either fixed margin leaves the claimed
$\epsilon N/4$ gap for all sufficiently large $N$ with probability
tending to one.  In the upper case, off-path bounds are strictly above
$\OPT$.  Before the optimum is found, best-bound search always has a
node on its path with bound at most $\OPT$; it cannot expand an off-path
node first.  Once the optimum is found, those nodes are pruned.  The
incumbent version follows directly from \Cref{lattice:path}.
\end{proof}

This is a threshold for a sufficient criterion for a short tree.  Failure
of C1 does not imply a superlinear tree lower bound.  The theorem gives no
conclusion at $\theta=\theta_c$ or in a shrinking critical window.
The simultaneous assertion follows from a finite coefficient tail; a
per-decoder limiting theorem such as
\citet{huLu2020LimitingPoissonMIMO} would not by itself supply a rate that
can be union-bounded over the $N$ wrong-fixing nodes.

\subsection{The logarithmic threshold for root box inactivity}
\label{bls:root-inactivity}

Let $u_N^\circ(\rho)$ be the unique minimizer of
\[
 \left\|w_N+\sqrt{\rho/N}\,H_Nu\right\|^2\quad\text{over }u\ge0,
\]
where the planted signs have again been absorbed into $H_N$.  Let
$R_N^\circ(\rho)$ denote this orthant value, and $R_N(\rho)$ the value
with $u\in[0,2]^N$.  We call equality of these values root box inactivity;
at the equality endpoint some coordinate can equal its upper bound.

\begin{theorem}[An exact random cutoff and its first-order scale]
\label{bls:root-threshold}
For every realization of full-rank $H_N$ and every $\rho>0$, define
\begin{equation}
 \rho_{\mathrm{box},N}=\frac14
           \left(\max_j u_{N,j}^\circ(1)\right)^2.
 \label{bls:random-root-cutoff}
\end{equation}
Then, pathwise,
\begin{equation}
 R_N(\rho)=R_N^\circ(\rho)
 \quad\Longleftrightarrow\quad \rho\ge\rho_{\mathrm{box},N}.
 \label{bls:pathwise-root-cutoff}
\end{equation}
If $\beta_N\to\beta\ge1$, then
\begin{equation}
 \frac{\rho_{\mathrm{box},N}}{\log N/(2\beta_N-1)}
       \xrightarrow{\Pbb}1.
 \label{bls:root-cutoff-scale}
\end{equation}
Consequently, for every fixed $\epsilon\in(0,1)$ and every deterministic
positive sequence $\rho_N$, root box inactivity holds with probability
tending to one if eventually
$\rho_N\ge(1+\epsilon)\log N/(2\beta_N-1)$.  If eventually
$\rho_N\le(1-\epsilon)\log N/(2\beta_N-1)$, then
$R_N(\rho_N)>R_N^\circ(\rho_N)$ with probability tending to one.
\end{theorem}

\begin{proof}
The substitution $z=\sqrt\rho\,u$ and uniqueness give
$u_N^\circ(\rho)=\rho^{-1/2}u_N^\circ(1)$.
The orthant minimizer lies in the box exactly when its maximum is at most
$2$.  Since the box is a subset of the orthant and the minimizer is unique,
this proves \eqref{bls:pathwise-root-cutoff}, including equality.

At $\rho=1$, \Cref{bls:nnls-law} gives the exact synthetic representation
\[
 \left(\max_j u_{N,j}^\circ(1)\right)^2
 \stackrel d=\frac NQ\max_{j\le K}|Z_j|^2,
 \quad K\sim\operatorname{Bin}(N,1/2),\quad
 Q\mid K\sim\chi^2_{M-K+1},
\]
with $Q$ and the standard normal coordinates independent conditional on
$K$.  The value is zero at $K=0$.  Binomial and chi-square concentration
give $K/N\to1/2$ and
$Q/N-(\beta_N-1/2)\to0$ in probability.  Elementary normal tail bounds
give $\max_{j\le K}|Z_j|^2/(2\log N)\to1$ in probability.
\Cref{app:bls-root-threshold-proof} proves these limits explicitly,
including the lower bound for the random number of coordinates.  Thus
\[
 \frac{\rho_{\mathrm{box},N}}{\log N/(2\beta_N-1)}
 =\frac{\beta_N-1/2}{Q/N}
   \frac{\max_{j\le K}|Z_j|^2}{2\log N}
 \xrightarrow{\Pbb}1.
\]
The two fixed-margin conclusions follow from the pathwise cutoff.
\end{proof}

A useful finite sufficient bound follows directly from
\Cref{bls:nnls-tail}:
\begin{equation}
 \Pbb\{\max_j u_{N,j}^\circ(\rho)>2\}
 \le\frac N2q^{M-N+1}\left(\frac{1+q}{2}\right)^{N-1},
 \qquad q=(1+4\rho/N)^{-1/2}.
 \label{bls:root-inactivity-tail}
\end{equation}
For $\rho=O(\log N)$, the logarithm of this bound is
$\log(N/2)-(2\beta_N-1)\rho+o(1)$.
The exact ratio $\beta_N$ is needed for this $o(1)$ remainder; substituting
its limit alone gives an $o(\log N)$ remainder under mere convergence.
The finite bound and the exact mixture both apply to square systems.
They describe equality with the orthant value.  They do not establish
planted recovery, successful rounding, a vertex root minimizer, or root
equality with the integer optimum.  The isolated value $\rho=0$ is excluded:
there the objective is constant in $u$ and the strict-convexity cutoff
argument does not apply.

\subsection{Certificate complexity below a linear signal level}
\label{bls:hard-certificates}

The lower bound is not obtained from C1 failure.  It uses large families
of flipped vertices and the value of their barycenter.  On coordinates
where the noise gives a descent direction, spreading a fixed number of
flips among many coordinates retains the linear improvement while
reducing the quadratic penalty.  A certifiable class must therefore have
substantial overlap among its flipped supports.  An entropy estimate
bounds the number of vertices that such a class can contain.

\begin{theorem}[Matching safe certificate regimes]
\label{bls:certificate-law}
Assume $\beta_N\to\beta\ge1$, and let
$0\le\varepsilon_N\le\rho_N$.  There are constants
$c_\beta>0$ and $\rho_0(\beta)>0$ with the following properties.
\begin{enumerate}[label=(\alph*)]
\item If $\rho_N\to\infty$ and $\rho_N=o(N)$, then with probability
tending to one,
\begin{equation}
 \log\kappa_{\OPT-\varepsilon_N}
 =\Theta_\beta\!\left(\frac N{\rho_N}\log\rho_N\right).
 \label{bls:sublinear-certificate-law}
\end{equation}
The minimum number of leaves of a variable-branching certificate has the
same logarithmic order.  A static variable order, fixed independently of
the data, attains the upper order with an optimal incumbent.
\item For each fixed $\rho\ge\rho_0(\beta)$, with probability tending to
one, both logarithms are $\Theta_{\beta,\rho}(N)$.
\end{enumerate}
In part (a), a static-order search given any incumbent value between
$\OPT$ and $W$ has the upper node bound for certifying that incumbent,
and an optimal incumbent makes it a certificate for the stated target.
\end{theorem}

\begin{proof}[Proof outline]
For the upper bound, take
$K=\lceil(1+o(1))N/[4(2\beta_N-1)\rho_N]\rceil$, with the $o(1)$
chosen positive and large enough to dominate the concentration errors.
Along a fixed variable order, simultaneously every node with $K$ wrong
fixings has bound at least $W$, with probability tending to one.  The
number of generated nodes is then at most
\begin{equation}
 1+2\sum_{i=1}^K\binom Ni
 \le1+2(eN/K)^K.
 \label{bls:static-tree-count}
\end{equation}
In the present sublinear regime $K\to\infty$, so its logarithm is at most
\begin{equation}
 (1+o(1))\frac{N}{4(2\beta-1)\rho_N}\log\rho_N.
 \label{bls:static-tree-exponent}
\end{equation}
An optimal incumbent can only prune earlier.  The full simultaneous
proof is in \Cref{app:bls-static-upper-proof}; it does not require planted
recovery.

For the lower bound, \Cref{app:bls-hard-lower-proof} proves
\begin{equation}
 \log\kappa_{\OPT-\varepsilon_N}
 \ge c_\beta\frac N{\rho_N}\log\rho_N-\log(8N)
 \label{bls:finite-hard-lower}
\end{equation}
in the divergent sublinear regime and for sufficiently large fixed
$\rho$.  It uses \eqref{bls:ml-value-gap} when $x^*$ has not been proved
optimal.  If $\rho_N\to\infty$ and $\rho_N=o(N)$, then
$\log N=o((N/\rho_N)\log\rho_N)$: for $\rho_N\le\sqrt N$ this is
immediate, and for $\rho_N>\sqrt N$ the ratio is at most
$2\rho_N/N$.  Thus \eqref{bls:finite-hard-lower} gives the lower order in
part (a).  The lower bound applies to all convex-piece certificates, so
also to variable trees.  At fixed sufficiently large $\rho$, it gives
$\Omega(N)$, while the full binary tree gives the upper $O(N)$ logarithm.
\end{proof}

For example, $\rho_N=c\log N$ with any fixed $c>0$ gives certificates of
size
\[
 \exp(\Theta(N\log\log N/\log N));
\]
planted recovery is justified
by \Cref{bls:ml} only when its constant condition is met.  At
$\rho_N=N^a$, $0<a<1$, the size is
$\exp(\Theta(N^{1-a}\log N))$.  A logarithmic root inactivity threshold
therefore coexists with a superpolynomial box certificate.

The constants in the lower proof are conservative.  In particular, the
subtractive $\log(8N)$ in \eqref{bls:finite-hard-lower} prevents extending
the matching statement to an arbitrary interval $\rho\le c'N$ without
checking constants.  At a fixed linear signal level the C1 theorem is
the positive result proved here; below its threshold, other short trees
are not excluded.  All these costs count certificates or nodes for the
specified oracle, and exclude the time needed to solve node relaxations.
